\section{Truncations divisible by the residue characteristic}
\label{sec:ramified-vertical-truncations}

A different integral lattice removes the remaining ramification
restriction for non-generation.  Its invariant lattice has covolume
\(n/p\), independently of the valuation of \(p\).
We use the whole-category primary index from
Corollary~\ref{cor:ramified-all-truncations}; the construction and
its exact-category compatibility are recalled below.

\begin{theorem}\label{thm:ramified-vertical-truncations}
Let \(O\) be a characteristic-zero DVR with residue characteristic
\(p>0\), let \(n\ge2\) be divisible by \(p\), and put
\(A=O[t]/(t^n)\), \(r=n/p\).
There is a class \(\xi\in K_0(\Lat_A(S_p))\) which vanishes both
under coefficient reduction \(t\mapsto0\) and after extension to
\(\operatorname{Frac}(O)\), and a homomorphism
\[
 \Theta:K_0(\Lat_A(S_p))\longrightarrow\Z/n
\]
which annihilates every Young permutation lattice and satisfies
\[
 \Theta(\xi)=r\pmod n.
\]
In particular, \(\Phi_A\) is not surjective in degree \(p\).
There is no restriction on the ramification or the residue field,
and \(O\) need not be complete.
\end{theorem}

Since \(0<r<n\), the displayed value is nonzero and has order \(p\).
Theorem~\ref{thm:vertical-ideal-exact-order} below proves that the
source class itself has exact order \(p\), and that its cyclic
subgroup splits from the whole Grothendieck group when \(p^2\nmid n\).

\begin{corollary}\label{cor:all-dvr-truncations-negative}
For every characteristic-zero DVR \(O\) of positive residue
characteristic and every \(n\ge2\), the ring \(O[t]/(t^n)\)
fails polynomial generation in degree equal to the residue
characteristic.
\end{corollary}

\begin{proof}
If \(p\mid n\), apply
Theorem~\ref{thm:ramified-vertical-truncations}.
If \(p\nmid n\), apply Theorem~\ref{thm:prime-to-p-all-dvr}.
\end{proof}

\subsection{The primary covolume on the whole exact category}

We first work over a complete \(O\), with uniformizer \(\pi\),
residue field \(k\), fraction field \(K\), and normalized valuation
\(v=v_\pi\).  Put \(P=C_p=\langle\sigma\rangle\), \(R=A[P]\), and
let \(\mathcal E=\Lat_A(P)\).
A map is stably zero if it factors through a finite free
\(R\)-module.

\begin{lemma}\label{lem:vertical-primary-index}
There is a homomorphism
\[
 \eta:K_0(\mathcal E)\longrightarrow\Z/n
\]
with the following properties.
It vanishes on every lattice obtained by coefficient extension
from \(O\).  If a lattice \(E\), with generator matrix \(T\), has a
lift \(\widetilde T\) over \(O[[t]]\) for which
\[
 U_E=\frac{\widetilde T^p-I}{t^n}\pmod{t^n}
\]
is the scalar \(-1\), then \(E\) has zero integral Tate cohomology
and
\[
 \eta(E)=v(\det E^P)\pmod n.
\]
The determinant is measured using an \(O\)-basis of \(E^P\) and
an \(A_K\)-basis of \((E_K)^P\), where \(A_K=K[t]/t^n\).
If instead \(U_E=-aI\), with \(a\in\{2,\ldots,p-1\}\), then
\(\eta(E)=0\).
\end{lemma}

\begin{proof}
This is the \(u+1\) coordinate of the construction in
Corollary~\ref{cor:ramified-all-truncations}.
We give its choice and exactness checks for the present setting.

For any \(E\in\mathcal E\), choose a lift \(\widetilde T\) and
define \(U_E\) by the displayed formula.
It is an \(R\)-endomorphism, since
\(\widetilde T^p-I\) commutes with \(\widetilde T\).
For an equivariant map \(F:E\to E'\), write
\[
 \widetilde T'\widetilde F-\widetilde F\widetilde T
       =t^n\widetilde B.
\]
Telescoping gives
\begin{equation}\label{eq:vertical-stable-curvature}
 U_{E'}F-FU_E
 =\sum_{i=0}^{p-1}T'^{\,p-1-i}BT^i
 =\sum_{j=0}^{p-1}T'^jHT^{-j},
 \qquad H=BT^{-1}.
\end{equation}
This group trace factors through \(R\otimes_A E\): use
\[
 v\longmapsto\sum_j\sigma^j\otimes T^{-j}v,\qquad
 \sigma^j\otimes w\longmapsto T'^jH(w).
\]
The same calculation compares lifts and bases.  Also the trace
of an equivariant map is \(p\) times that map.  Thus stable Hom
groups are killed by \(p\), and \(U_E\) is stably natural and
independent of those choices.

Let \(E_-\) be the integral Fitting summand on which \(U_E+I\)
is nilpotent modulo \(\pi\).  Its complement has \(U_E+I\)
invertible.  The Fitting decomposition modulo \(\pi\) lifts
uniquely over the complete DVR by compatible Chinese-remainder
projectors on \(E/\pi^qE\).  These projectors are polynomials in
\(U_E\), commute with the \(A\)- and \(P\)-actions, and give
\(A\)-free summands, since \(A\) is local.  For a sequence with
intertwined endomorphisms, a common annihilating polynomial gives
compatible projectors and an exact sequence of these summands.
This uses only a linear primary polynomial, so it requires no
finiteness or perfection assumption on \(k\).

The resulting summand is independent of the lift up to stable
isomorphism.  Write \(p=u\pi^e\) with \(u\in O^\times\).
On \(E_-\), some power of \(U_E+I\) is divisible by \(\pi\).
Its \(e\)-th power is consequently divisible by \(p\), so is
stably zero.  Division by \(\pi\) or \(p\) preserves equivariance
because the lattices have no such torsion.
On the complementary summand \(U_E+I\) is invertible.
Stable naturality in \eqref{eq:vertical-stable-curvature}
therefore makes every stable map between a selected summand and
an unselected summand zero.  The identity comparing two choices,
and its inverse, are consequently block diagonal in the stable
category and induce stable inverse maps on the selected summands.

The curvature on \(E_-\) is invertible, since its residue is
primary at \(-1\ne0\).  We check integral Tate acyclicity.
Choose an \(A\)-basis adapted to the summand and lift it to
\(O[[t]]\).  Off-diagonal blocks of the conjugated generator
lift are divisible by \(t^n\).  Every contribution leaving a
diagonal block and returning is divisible by \(t^{2n}\).
Thus its diagonal block \(\widetilde T_-\) has curvature
congruent to \(U_E|_{E_-}\) modulo \(t^n\), and this curvature
\(\widetilde U_-\) is a unit over \(O[[t]]\).
Write \(\operatorname{Nm}(Z)=I+Z+\cdots+Z^{p-1}\).
If \((T_--I)x=0\), lift \(x\) and write
\((\widetilde T_--I)x=t^ny\).  Multiplication by
\(\operatorname{Nm}(\widetilde T_-)\) and cancellation of \(t^n\)
give
\[
 \widetilde U_-x=\operatorname{Nm}(\widetilde T_-)y.
\]
Since \(\widetilde U_-\) commutes with the lift, \(x\) is a norm
boundary modulo \(t^n\).
If instead \(\operatorname{Nm}(T_-)x=0\), write
\(\operatorname{Nm}(\widetilde T_-)x=t^ny\) and multiply by
\(\widetilde T_--I\).  This gives
\[
 x=(\widetilde T_--I)\widetilde U_-^{-1}y
 \pmod{t^n}.
\]
Both Tate groups therefore vanish.

For any cohomologically trivial coefficient-free lattice \(V\),
the invariant module \(V^P\) is \(O\)-free.  The projector
\(\operatorname{Nm}(T)/p\) makes \((V_K)^P\) projective, hence
free, over the local ring \(A_K\).  Define its covolume
\[
 \iota(V)=v(\det V^P)\pmod n
\]
using an \(A_K\)-basis and the associated \(K\)-basis
\(1,t,\ldots,t^{n-1}\) in each coordinate.
An \(O\)-basis change has determinant of valuation zero.
For \(B\in\operatorname{GL}_d(A_K)\), the induced map on each
successive \(t\)-adic quotient is its constant matrix \(B_0\),
so
\[
 \det_K(B)=\det(B_0)^n.
\]
Thus the generic basis ambiguity is a multiple of \(n\).

The covolume is additive on admissible sequences of
cohomologically trivial lattices.  Their invariant modules
form a short exact sequence because \(H^1(P,V')=0\) for the
left term.  Compatible generic bases and \(O\)-bases give a
block triangular determinant.  It vanishes on a free \(R\)-module,
whose invariants have group-norm elements as an \(A\)-basis.
It is also invariant under stable isomorphism.
Indeed, for stable inverse maps \(f:V\to W\) and \(g:W\to V\),
write \(gf-I=ba\) through a free module \(Q_0\).
Then \((f,a):V\to W\oplus Q_0\) is split by \((g,-b)\).
The complementary projection is
\[
 \begin{pmatrix}I-fg&fb\\-ag&I+ab\end{pmatrix}.
\]
Every entry is stably zero; hence its image is projective.
The ring \(R\) is local, since its residue algebra is
\(k[P]\simeq k[X]/X^p\).  The complement is therefore free,
so the two lattices become isomorphic after adjoining free
modules, whose covolumes vanish.

Set \(\eta(E)=\iota(E_-)\).
Stable independence of \(E_-\) and the preceding argument prove
choice independence.  For an admissible exact sequence, choose
an \(A\)-splitting and lift the generator in upper triangular
blocks.  The curvature is upper triangular with the endpoint
curvatures on its diagonal.  The compatible Fitting projectors
give an admissible sequence
\[
 0\longrightarrow E'_-\longrightarrow E_-
       \longrightarrow E''_-\longrightarrow0.
\]
All three terms are cohomologically trivial, so their covolumes
add.  Choice independence permits these compatible lifts for
every original sequence.  This proves the homomorphism on the
whole \(K_0(\mathcal E)\).
For a constant \(O\)-lattice, a constant lift has curvature zero,
so the selected summand is zero.
Finally, \(U_E=-I\) selects all of \(E\), whereas
\(U_E=-aI\), \(a\ne1\) in \(\F_p^\times\), selects none.
\end{proof}

\subsection{A finite ideal lattice over every DVR}

The witness construction is defined over \(O\) before completion.
Choose a uniformizer \(\pi\), set \(n=pr\), and form
\[
 C=O[t],\qquad
 F(Y,t)=(1+Y-t^r)^p-1+t^n,\qquad
 S=C[Y]/(F).
\]
Write
\begin{equation}\label{eq:vertical-polynomial-coefficients}
 \begin{split}
 F(Y,t)&=Y^p+\sum_{j=0}^{p-1}q_j(t)Y^j,\\
 q_0(t)&=(1-t^r)^p-1+t^n,\\
 q_j(t)&=\binom pj(1-t^r)^{p-j}\qquad(1\le j<p).
 \end{split}
\end{equation}
Every \(q_j(t)\) is divisible by \(p\), hence by \(\pi\).
For \(q_0\), this follows either by expansion or from
\((1-t^r)^p\equiv1-t^{pr}\pmod p\).
In particular \(F(Y,t)\equiv Y^p\pmod\pi\).
Put
\[
 N=(\pi,Y)\subset S.
\]

\begin{lemma}\label{lem:vertical-ideal-lattice}
The \(C\)-module \(N\) is free with basis
\[
 e_0=\pi,\qquad e_j=Y^j\quad(1\le j\le p-1).
\]
The module \(M=N/t^nN\) is a free \(A\)-module of rank \(p\).
Multiplication by \(1+Y-t^r\) gives a \(P\)-action on \(M\)
and has an integral polynomial lift \(\widetilde T\) satisfying
\begin{equation}\label{eq:vertical-exact-curvature}
 \widetilde T^p=I-t^nI.
\end{equation}
Under \(X=Y-t^r\), there is an equivariant inclusion
\begin{equation}\label{eq:vertical-ideal-inclusion}
 M\hookrightarrow R=A[X]/(g(X))
\end{equation}
with image \((\pi,X+t^r)\), and its quotient is \(k[t]/t^n\).
\end{lemma}

\begin{proof}
Since \(F\) is monic, \(S\) is \(C\)-free on the powers
\(1,Y,\ldots,Y^{p-1}\).
The span of \(\pi,Y,\ldots,Y^{p-1}\) is stable under \(Y\):
the only relation to check is \(Y^p=-\sum_jq_jY^j\), whose
constant term is divisible by \(\pi\).
This span is therefore precisely the ideal \((\pi,Y)\).
Its generators are \(C\)-independent because \(\pi\) is a
nonzerodivisor on \(C\).
This proves both freeness assertions.

The defining equation of \(S\) gives
\((1+Y-t^r)^p=1-t^n\).
It acts on the \(C\)-free module \(N\), proving
\eqref{eq:vertical-exact-curvature}; modulo \(t^n\) this is a
group action.  Explicitly, in the indicated basis,
\[
 \begin{split}
 \widetilde T e_0&=(1-t^r)e_0+\pi e_1,\\
 \widetilde T e_j&=(1-t^r)e_j+e_{j+1}
                       \qquad(1\le j\le p-2),\\
 \widetilde T e_{p-1}
   &=-\frac{q_0}{\pi}e_0
     -\sum_{j=1}^{p-2}q_j e_j
     +(1-t^r-q_{p-1})e_{p-1}.
 \end{split}
\]
The sums indexed from \(1\) to \(p-2\) are empty at \(p=2\).
Every displayed coefficient lies in \(O[t]\).

Since \(S/N=k[t]\), multiplication by \(t^n\) is injective
on \(S/N\).  Tensoring \(0\to N\to S\to S/N\to0\) with
\(C/(t^n)\) therefore preserves injectivity on the left.
The substitution \(X=Y-t^r\) identifies \(S/t^nS\) with
\(A[X]/(g(X))=R\).
It gives \eqref{eq:vertical-ideal-inclusion}, the stated ideal
image, and the quotient \(k[t]/t^n\).
\end{proof}

\subsection{The invariant lattice and the two specialization kernels}

\begin{lemma}\label{lem:vertical-invariant-volume}
For the lattice \(M\) of Lemma~\ref{lem:vertical-ideal-lattice},
\[
 M^P=(\pi,t^r)h(X)\subset A_Kh(X)=(M_K)^P.
\]
Consequently, over complete \(O\),
\begin{equation}\label{eq:vertical-index-value}
 \eta(M)=r\pmod n.
\end{equation}
The cyclic class \(x=[M]-[R]\) vanishes under \(t\mapsto0\)
and over \(K=\operatorname{Frac}(O)\).
\end{lemma}

\begin{proof}
In the regular module \(R=A[X]/(Xh)\), the kernel of \(X\)
is exactly \(Ah(X)\).  For example, divide a representative
by the monic polynomial \(Xh\) to obtain degree less than \(p\).
If its product with \(X\) is zero, that product is a scalar
multiple of \(Xh\); cancellation of \(X\) in \(A[X]\) gives
a scalar multiple of \(h\).  Its coefficient of \(X^{p-1}\)
also proves that \(A\to Ah\) is injective.

Using the inclusion \eqref{eq:vertical-ideal-inclusion},
membership of \(a h(X)\) in \(M\) is tested in
\[
 R/M=k[t]/t^n,\qquad X=-t^r.
\]
Modulo \(\pi\), one has \(h(X)=X^{p-1}\), so its image is
\((-1)^{p-1}t^{r(p-1)}=(-1)^{p-1}t^{n-r}\).
Its annihilator in \(k[t]/t^n\) is \((t^r)\).
Thus \(a h(X)\) belongs to \(M\) exactly when
\(a\in(\pi,t^r)\), proving the formula for the invariants.

After inverting \(\pi\), the ideal \(M\subset R\) becomes all
of \(R_K\); in particular \(h(X)\) is an \(A_K\)-basis for the
invariant module.  Relative to that basis, the quotient of
\(A\) by the invariant lattice is
\[
 A/(\pi,t^r)\simeq k[t]/t^r,
\]
of \(O\)-length \(r\).
Equation~\eqref{eq:vertical-exact-curvature} gives the actual
scalar curvature \(-I\), so
Lemma~\ref{lem:vertical-primary-index} applies and gives
\eqref{eq:vertical-index-value}.

The same generic ideal equality proves \(M_K\simeq R_K\).
For reduction, multiplication by \(t\) is also regular on
\(S/N=k[t]\).  Therefore \(M/tM=N/tN\) injects into
\(O[P]\) as the ideal
\[
 I_\pi=(\pi,X).
\]
Let \(I=XO[P]\) be the augmentation ideal.
The augmentation sequences
\[
 0\longrightarrow I\longrightarrow I_\pi
        \longrightarrow \pi O\longrightarrow0,\qquad
 0\longrightarrow I\longrightarrow O[P]
        \longrightarrow O\longrightarrow0
\]
split over \(O\); the quotients are trivial lattices and
\(\pi O\simeq O\).
Hence \([I_\pi]=[O[P]]\) in \(K_0(\Lat_O(P))\).
This proves the reduction assertion, without asserting that
the two reduced lattices are isomorphic.
\end{proof}

\subsection{Normalizer twists, induction, and descent}

\begin{proof}[Proof of Theorem~\ref{thm:ramified-vertical-truncations}]
First suppose \(O\) is complete.
For \(1\le a\le p-1\), let \(M^{(a)}\) be the twist on which
\(\sigma\) acts by \(T^a\).  Use \(\widetilde T^a\) as a lift.
Equation~\eqref{eq:vertical-exact-curvature} gives
\[
 \frac{\widetilde T^{ap}-I}{t^n}
 =\frac{(1-t^n)^a-1}{t^n}I
 \equiv -aI\pmod{t^n}.
\]
The invariants of \(M^{(a)}\) are the same as those of \(M\),
because \(\sigma^a\) generates \(P\).  Thus
Lemma~\ref{lem:vertical-primary-index} gives
\begin{equation}\label{eq:vertical-normalizer-values}
 \eta(M^{(a)})=
 \begin{cases}
 r\pmod n,&a=1,\\
 0,&a\ne1.
 \end{cases}
\end{equation}
Every twist of \(R\) is a constant lattice, so its value is zero.

Embed \(P\) as a \(p\)-cycle in \(S_p\), and put
\[
 x=[M]-[R],\qquad
 \xi=\operatorname{Ind}_{P}^{S_p}x
     =[\operatorname{Ind}_{P}^{S_p}M]-[A[S_p]],
 \qquad
 \Theta=\eta\circ\operatorname{Res}_{P}^{S_p}.
\]
For a double coset outside the normalizer of \(P\), the
intersection of \(P\) with its conjugate is trivial, since
\(P\) has prime order.
The corresponding Mackey term factors through
\(\operatorname{Res}_{1}^{P}x=0\): both \(M\) and \(R\)
are coefficient-free of rank \(p\).
Inside the normalizer, \(N_{S_p}(P)/P\simeq\F_p^\times\)
gives one copy of every automorphism twist.
Equation~\eqref{eq:vertical-normalizer-values} therefore yields
\[
 \Theta(\xi)=
 \sum_{a\in\F_p^\times}\eta(x^{(a)})
 =r\pmod n.
\]
Changing the convention for conjugate modules replaces the
labels by their inverses, which permutes the same set and
leaves its unique identity contribution unchanged.

Every Young permutation lattice is obtained from an integral
\(O\)-permutation lattice by coefficient extension to \(A\).
The same is true after restriction to \(P\).
Lemma~\ref{lem:vertical-primary-index} therefore gives
\(\Theta=0\) on the Young span.
Restriction and induction preserve coefficient freeness and
coefficient-split exact sequences, so all these maps are maps
on the specified exact Grothendieck groups.
The two kernel statements for \(x\), proved in
Lemma~\ref{lem:vertical-invariant-volume}, commute with
induction and give the assertions for \(\xi\).

For an arbitrary DVR \(O\), construct \(N\), \(M\), and \(\xi\)
by the same finite formulas over \(O[t]\).
The bases, ideal inclusions, augmentation sequences, and
generic equality above do not use completeness.
Extend coefficients to the completion \(\widehat O\).
This is flat, preserves coefficient-split sequences and
coefficient-free objects, and retains the uniformizer,
residue field, and normalized valuations.
Compose the completed \(\Theta\) with this scalar-extension
map.  The invariant quotient remains \(k[t]/t^r\), so the
value on the original \(\xi\) is still \(r\pmod n\).
The Young span is still annihilated.
This proves the theorem without a completeness assumption.
\end{proof}

\begin{example}\label{ex:vertical-pure-power-tests}
The new ideal lattice differs from the standard
\(g+t^n\) mesh of Section~\ref{sec:pure-power-vanishing}.
For \(p=n=3\), \(O=\Z_3[\pi]\), \(\pi^3=3\), its matrix modulo
\(t^3\), in the basis \(\pi,Y,Y^2\), is
\[
 T=
 \begin{pmatrix}
 1-t&0&\pi^2(t-t^2)\\
 \pi&1-t&-3(1-t)^2\\
 0&1&-2(1-t)
 \end{pmatrix}.
\]
The invariant ideal is \((\pi,t)\), and the selected coordinate
is \(1\) modulo \(3\).
For \(p=2,n=4\), \(O=\Z_2[\pi]\), \(\pi^4=2\), a matrix is
\[
 T=
 \begin{pmatrix}
 1-t^2&\pi^3t^2\\
 \pi&-1+t^2
 \end{pmatrix}
 \quad\text{over }O[t]/t^4.
\]
Here the invariant ideal is \((\pi,t^2)\) and the coordinate is
\(2\) modulo \(4\).
For the exact polynomial lift in
\eqref{eq:vertical-exact-curvature}, its upper-right entry is
\(\pi^3(t^2-t^4)\).
These nonzero classes coexist with the vanishing standard mesh
classes proved in Theorem~\ref{thm:pure-power-mesh-vanishing}; they are different integral lattices.
\end{example}
